\section{Local scarcity at matched scales}
\label{loc:section}

This section proves \cref{can:local-scarcity}. We use its two outward
paths and the notation established in the cancellation argument. In
particular, $c=99/100$, the central coordinate is fixed, and the variable
is $\theta=q_{e-1}$. The selected indices are denoted by $j_+$ and
$j_-$. Every assertion concerning an interval in $\theta$ uses real
lifts of the recurrence. The constants below are independent of the
selected indices and of the base configuration.

The central relation in \cref{can:center-equation} will become, after rescaling, an approximate
identity between a cubic polynomial and two shifted reciprocal-cosine
terms. A positive-measure Hausdorff limit of the closures of successful
sets would force an analytic identity. We rule out that identity even when the two
reciprocal terms nearly cancel or one phase slope tends to zero. The
derivative estimates in the first two subsections provide the uniform
remainders needed for these two degeneracies.

\subsection{Control on a single rescaled path}

Fix either outward path and write $j$ for its selected index. Put
\[
 L_j=\abs{t_j(\theta_*)},\qquad
 \xi=L_j(\theta-\theta_*),\qquad D=\frac{\dd}{\dd\xi}.
\]
The initial position $q_0$ is constant and $q_1$ is affine in $\xi$,
with derivative $\pm L_j^{-1}$. By \cref{can:good-t}, the base
linearized solution satisfies
\begin{equation}
 \frac{\abs{t_l(\theta_*)}}{L_j}
 \le M^{-c(j-l)}\quad(1\le l\le j).
 \label{loc:base-growth}
\end{equation}
We will prove, throughout $\abs{\xi}\le1$, that
\begin{align}
 \abs{q_l(\xi)-q_l(0)}+\abs{Dq_l(\xi)}
 &\le C M^{-c(j-l)}, \label{loc:first-derivatives}\\
 \abs{D^p q_l(\xi)}
 &\le C_p M^{-pc(j-l)}M^{1-pc}
 \quad(2\le p\le5), \label{loc:higher-derivatives}
\end{align}
for $0\le l\le j$. The terms with $l=0$ vanish.

Write $B_l=M^{-c(j-l)}$ and $\Delta_l=q_l(\xi)-q_l(0)$.
Taylor's formula for the recurrence, followed by propagation with the
base linearized matrices, gives
\begin{equation}
 \abs{\Delta_l}
 \le \abs{\xi}\frac{\abs{t_l(\theta_*)}}{L_j}
       +C\sum_{d=1}^{l-1}M^{l-d}\abs{\Delta_d}^2.
 \label{loc:nonlinear-comparison}
\end{equation}
Here the quadratic error at site $d$ is at most $CM\abs{\Delta_d}^2$,
and the remaining $l-d-1$ steps have norm at most $M^{l-d-1}$.
For $l=1$ the estimate $\abs{\Delta_l}\le2B_l$ follows directly
from \cref{loc:base-growth}. If it holds for all earlier indices,
the sum in \cref{loc:nonlinear-comparison} is bounded by
\[
 4C B_l^2\sum_{a\ge1}M^{(1-2c)a}.
\]
Since $B_l\le1$ and $1-2c<0$, this is at most $B_l$ for all
sufficiently large $M$. Induction proves the required displacement
bound, uniformly in $j$.

Compare $Dq_l(\xi)$ with the base variation
$\pm t_l(\theta_*)/L_j$. Their initial data agree. Their difference
is propagated by the base matrices with forcing
\[
 \bigl(v(q_d(\xi))-v(q_d(0))\bigr)Dq_d(\xi),
\]
whose absolute value is at most $CM\abs{\Delta_d}\abs{Dq_d}$.
Induction with $\abs{Dq_d}\le2B_d$ therefore uses the same geometric
sum and proves \cref{loc:first-derivatives}.

For completeness, the higher derivative induction uses the recurrence
at the varied orbit. For $p\ge2$ it has the form
\[
 D^p q_{l+1}=v(q_l)D^p q_l-D^p q_{l-1}+F_{p,l},
\]
where $F_{p,l}$ is a finite sum, with coefficients depending only on
$p$, of terms
\[
 \phi^{(s)}(q_l)\prod_{i=1}^{s}D^{m_i}q_l,
 \qquad 2\le s\le p,\quad m_i\ge1,\quad
 \sum_{i=1}^{s}m_i=p.
\]
Each $m_i$ is strictly smaller than $p$. The first derivative bound
and the induction hypothesis give
$\abs{D^{m_i}q_l}\le C_{m_i}B_l^{m_i}$, since
$M^{1-m_ic}\le1$ whenever $m_i\ge2$. Consequently
$\abs{F_{p,l}}\le C_p M B_l^p$. The initial higher derivatives
at sites $0,1$ vanish. Propagation by the varied matrices gives
\[
 \abs{D^p q_l}
 \le C_p\sum_{d=1}^{l-1}M^{l-d}B_d^p
 \le C_p B_l^p\sum_{a\ge1}M^{(1-pc)a}
 \le C'_p B_l^p M^{1-pc}.
\]
This proves \cref{loc:higher-derivatives} successively for
$p=2,3,4,5$.

At the selected site $Dq_j(0)=\pm1$. The case $p=2$ shows that
\[
 Dq_j(\xi)=Dq_j(0)+O(M^{1-2c}).
\]
It follows, for all sufficiently large $M$, that
\begin{equation}
 \abs{t_j(\xi)}\ge L_j/2,\qquad
 \abs{t_l(\xi)}\le C L_j M^{-c(j-l)}\quad(l<j).
 \label{loc:varied-t}
\end{equation}
The quotient
\begin{equation}
 E=\frac{t_{j-1}}{t_j}=\frac{Dq_{j-1}}{Dq_j}
 \quad\hbox{satisfies}\quad
 \abs{D^p E}\le C_p M^{-c}\quad(0\le p\le4).
 \label{loc:quotient}
\end{equation}
Indeed the denominator is bounded away from zero. Every derivative of
its reciprocal of order at most four is bounded by the already proved
bounds through $D^5q_j$. Every derivative of the numerator through
order four contains a derivative of $q_{j-1}$ and has a factor at most
$CM^{-c}$. The product rule proves \cref{loc:quotient}.

\subsection{A fourth derivative bound for the boundary value}

Fourth-order control makes the boundary value replaceable by a cubic
polynomial with a small remainder. The remainder must remain negligible
even after normalizing the first nonzero correction when the two
reciprocal-cosine terms nearly cancel.

Let $R_j$ be the initial value at site $1$ of the solution $U$ with
$U_0=1$ and $U_j=0$. This solution exists throughout our interval
because $t_j$ does not vanish. We claim that, with $\rho=4/5$,
\begin{equation}
 L_j^2\abs{D^4R_j}\le C M^{-2-\rho}.
 \label{loc:boundary-fourth}
\end{equation}

Consider the $(j-1)$-by-$(j-1)$ tridiagonal matrix $B$ with diagonal
$v(q_a)$, $1\le a\le j-1$, and both adjacent diagonals equal to
$-1$. Its determinant is $t_j$: expansion along the last row gives
the same recursion and initial conditions as those for $t$.
The interior boundary problem is $BU=e_1$, so $R_j=(B^{-1})_{11}$.
More explicitly, its inverse $G=B^{-1}$ is
\begin{equation}
 G_{ab}=t_{\min(a,b)}U_{\max(a,b)}.
 \label{loc:green}
\end{equation}
To check this formula, fix column $b$. Its entries solve the homogeneous
recurrence except at row $b$, and are zero at the boundary sites $0,j$.
At row $b$ the output of $B$ is
$t_{b+1}U_b-t_bU_{b+1}=1$, by the constant Wronskian and the data
$t_0=0,t_1=1,U_0=1$. Thus the column is precisely $B^{-1}e_b$.

At the right boundary the same Wronskian gives
$U_{j-1}=1/t_j$. Backward propagation and
\cref{loc:varied-t} yield
\begin{equation}
 L_j\abs{U_a}\le2M^{j-a-1}.
 \label{loc:boundary-backwards}
\end{equation}
For $d_a=j-a\ge1$, \cref{loc:green,loc:varied-t,loc:boundary-backwards}
therefore give
\begin{equation}
 \abs{G_{ab}}\le C
 M^{\min(d_a,d_b)-1-c\max(d_a,d_b)}.
 \label{loc:green-size}
\end{equation}
The chain rule and the bounds for derivatives of $q_a$ also give
\begin{equation}
 \abs{D^p v(q_a)}\le C_p M^{1-cp d_a}
 \quad(1\le p\le4).
 \label{loc:potential-derivatives}
\end{equation}
In fact every term contains a derivative of $v$ bounded by $C_pM$
and derivatives of $q_a$ whose total order is $p$.

Starting with $DG=-G(DB)G$, differentiate three more times and take
the $(1,1)$ entry. The result is a finite linear combination, with
fixed integer coefficients, of sums of terms
\begin{equation}
 U_{a_1}(D^{p_1}v_{a_1})G_{a_1a_2}
 (D^{p_2}v_{a_2})\cdots
 G_{a_{s-1}a_s}(D^{p_s}v_{a_s})U_{a_s},
 \qquad p_i\ge1,\quad \sum_{i=1}^{s}p_i=4.
 \label{loc:inverse-chains}
\end{equation}
Here $v_a=v(q_a)$, $1\le s\le4$, and each site index is summed
independently from $1$ to $j-1$. This form follows inductively:
differentiating a potential factor increases its derivative order,
whereas differentiating any inverse factor inserts another
$-G(DB)G$. Also $G_{1a}=U_a$ because $t_1=1$.

Put $d_i=j-a_i$. After multiplication by $L_j^2$, the exponent of
$M$ bounding a term of \cref{loc:inverse-chains} is
\begin{align*}
 H&=-1+d_1+d_s+
 \sum_{i=1}^{s-1}\bigl(\min(d_i,d_{i+1})-c\max(d_i,d_{i+1})\bigr)
       -c\sum_{i=1}^{s}p_id_i\\
 &\le -1+2\min_i d_i-(3c-2)\sum_{i=1}^{s}p_id_i\\
 &\le -1-\frac{47}{100}\sum_{i=1}^{s}p_id_i.
\end{align*}
For the first inequality, write each adjacency term as
$-\abs{d_i-d_{i+1}}+(1-c)\max(d_i,d_{i+1})$. The total variation
is at least $d_1+d_s-2\min_i d_i$, while the sum of maxima is at
most $2\sum_i d_i\le2\sum_i p_id_i$. For the last inequality use
$4\min_i d_i\le\sum_i p_id_i$ and $c=99/100$.
Summing over sites is harmless: for every composition of $4$,
\[
 \sum_{d_1,\ldots,d_s\ge1}
 M^{-1-(47/100)\sum_i p_id_i}
 =M^{-1}\prod_{i=1}^{s}
 \frac{M^{-(47/100)p_i}}{1-M^{-(47/100)p_i}}
 \le C M^{-72/25}.
\]
There are only finitely many such compositions and chain types.
Since $72/25>2+4/5$, this proves \cref{loc:boundary-fourth}.

\subsection{The cubic relation at successful points}

Return to the two paths and put
$L=\max(L_{j_+},L_{j_-})$, $h=L(\theta-\theta_*)$.
Number the paths so that the first has the larger base scale. Write
$Q_i(h)$ for their selected positions and $E_i(h)$ for the quotients
in \cref{loc:quotient}. With $b_i=Q_i'(0)$ and $d_i=Q_i''(0)$,
the preceding estimates imply
\begin{align}
 Q_i(h)&=z_i+b_ih+\tfrac12d_ih^2+\mathcal E_i(h),
 &\norm{\mathcal E_i}_{C^1([-1,1])}
   &\le C\abs{b_i}^3M^{1-3c}, \label{loc:phase-taylor}\\
 \abs{d_i}&\le C\abs{b_i}^2M^{1-2c},
 &\abs{b_1}&=1,\quad
 \abs{b_2}=r\in[M^{-3/5},1], \label{loc:phase-slopes}\\
 \abs{\partial_h^pE_i}&\le C_pM^{-c}\abs{b_i}^p
 &&(0\le p\le4). \label{loc:scaled-quotient}
\end{align}
The scale change is $\xi=(L_{j_i}/L)h=\abs{b_i}h$.
It also gives, for $2\le p\le5$,
\begin{equation}
 \abs{Q_i^{(p)}}\le C_p\abs{b_i}^pM^{1-pc},\qquad
 \abs{Q_i'}\asymp\abs{b_i}.
 \label{loc:scaled-phase-derivatives}
\end{equation}

At a successful point there is a common solution with central value
$1$, and its two outward values at site $1$ add to $v(q_e)$.
Apply \cref{can:boundary-expansion} to both paths at that point.
Since $\abs{t_{j_i}(\theta)}=L\abs{Q_i'(h)}$ and
$v(Q_i)-E_i=K(\cos(2\pi Q_i)+(2-E_i)/K)$, multiplication by
$KL^2$ yields
\[
 KL^2\left(v(q_e)-\sum_{i=1}^{2}R_{j_i}(\theta)\right)
 =\sum_{i=1}^{2}
 \frac{1}{(Q_i')^2(\cos(2\pi Q_i)+(2-E_i)/K)}
 +O(r^{-2}M^{1-3c}).
\]
The error bound uses \cref{loc:scaled-phase-derivatives} to compare
the varied $t_{j_i}$ with their base sizes. Replace the left side by
its cubic Taylor polynomial $P(h)$ at $0$. The contribution of each
$R_{j_i}$ to the Taylor remainder is at most
\[
 C K L^2\abs{b_i}^4 L_{j_i}^{-2}M^{-2-\rho}
 =C K\abs{b_i}^2M^{-2-\rho}.
\]
Consequently every successful $h\in[-1,1]$ satisfies
\begin{equation}
 P(h)=\sum_{i=1}^{2}
 \frac{1}{(Q_i'(h))^2
       (\cos(2\pi Q_i(h))+(2-E_i(h))/K)}
       +O(M^{-1-\rho}+r^{-2}M^{1-3c}).
 \label{loc:cubic-relation}
\end{equation}
The polynomial may depend on the configuration. Its degree is at
most three. The displayed error is uniform on the successful set,
and tends to zero uniformly over all configurations, since
$r^{-2}M^{1-3c}\le M^{-77/100}$.

\subsection{Polynomial extraction from a Hausdorff limit}

Consider any sequence of configurations with $M\to\infty$ and
nonempty successful sets $\mathcal S_M$. Pass to a subsequence on
which their closures converge in Hausdorff distance to a nonempty
compact set $F\subset[-1,1]$. We prove that $F$ has measure zero.
Suppose to the contrary that it has positive measure.

We record precisely the compactness argument for the polynomials
that will be used below. Suppose, after an indicated subtraction
or multiplication of \cref{loc:cubic-relation}, a polynomial $P_M^*$
of degree at most three agrees at successful points, up to an error
tending to zero, with functions converging uniformly on every compact
subinterval of $[-1,1]$ away from a fixed finite set of limiting poles
in that interval. Choose four
distinct points of $F$ outside those poles. Hausdorff convergence
and approximation of points in the closure supply successful points
converging to these four points. The values of $P_M^*$ there are
bounded. The associated Vandermonde matrices converge to an invertible
matrix, so all four coefficients of $P_M^*$ are bounded. Pass to a
further subsequence on which the polynomials converge coefficientwise.
For every point of $F$ away from the poles the same approximation
argument shows equality with the limiting function. In particular,
equality holds on a set of positive measure. Clearing the finitely
many denominator factors used in that case gives an identity between
real analytic functions on $\R$: the difference has an accumulation
point of zeros and hence vanishes identically. Passing to further
subsequences throughout does not change the fixed Hausdorff limit $F$.

\subsection{Nonzero limiting slopes and opposite cosines}

First suppose that $r$ has a positive limit $r_0$. Take convergent
subsequences of the phases modulo integers and of the slopes. The
first slope has a fixed sign after a further subsequence, so write
its value as $b\in\{-1,1\}$. By
\cref{loc:phase-taylor,loc:scaled-quotient}, the two terms on the
right of \cref{loc:cubic-relation} converge, away from their limiting
poles, to
\[
 \sec(2\pi(z+bh))+
 r_0^{-2}\sec(2\pi(z'+b'h)),\qquad \abs{b'}=r_0.
\]
Polynomial extraction and multiplication by both cosines show that
the zero sets of these cosines on $\R$ coincide: at a zero of either
cosine, the other must vanish because its coefficient is positive.
Their zero sets are arithmetic progressions with respective spacings
$1/(2\abs{b})$ and $1/(2\abs{b'})$, so $r_0=1$.
The two cosines are therefore equal or opposite. Equality would leave
twice a secant, which cannot equal a polynomial, as multiplication by
the cosine and evaluation at one of its zeros shows. Only the opposite
case remains.

Reflect the second phase if necessary and translate it by an odd
half-integer. The resulting phase $\widetilde Q_2$ has the same
limiting affine phase as $Q_1$, while
\[
 \cos(2\pi Q_2)=-\cos(2\pi\widetilde Q_2),\qquad
 (Q_2')^2=(\widetilde Q_2')^2.
\]
By \cref{loc:phase-taylor},
\begin{equation}
 \widetilde Q_2-Q_1=p(h)+O_{C^1}(M^{1-3c}),\qquad
 \deg p\le2,\qquad \norm{p}_{\rm coeff}\longrightarrow0.
 \label{loc:phase-difference}
\end{equation}
Here $\norm{p}_{\rm coeff}$ is the maximum absolute value of its
three coefficients. Define
\[
 \ell=\max(\norm{p}_{\rm coeff},K^{-1}).
\]
Then $\ell\to0$, but every error of order $M^{1-3c}=M^{-197/100}$
is $o(\ell)$. The error in \cref{loc:cubic-relation} is also
$o(\ell)$ in the present case, since $r\to1$ and $\rho>0$.

We expand on a compact subinterval away from the zeros of the common
limiting cosine. Put
\[
 f(x,y)=y^{-2}\sec(2\pi x),\qquad
 A=2\pi(z_1+b_1h).
\]
The denominators in \cref{loc:cubic-relation} stay uniformly away
from zero there for large $M$. Their shifts are
$2/K+O(M^{-c}/K)$. Expanding the reciprocals, with their respective
cosine signs, gives the sum of the two terms as
\begin{align*}
 &f(Q_1,Q_1')-f(\widetilde Q_2,\widetilde Q_2')\\
 &\quad-\frac2K\left((Q_1')^{-2}\sec^2(2\pi Q_1)
            +(\widetilde Q_2')^{-2}\sec^2(2\pi\widetilde Q_2)\right)
       +o(\ell).
\end{align*}
The reciprocal expansion has error $O(K^{-2}+M^{-c}/K)=o(\ell)$.
It is essential to subtract the two $f$ terms at the actual common
phase before replacing that phase by an affine one. Taylor expansion
of $f$, using \cref{loc:phase-difference}, gives
\[
 f(Q_1,Q_1')-f(\widetilde Q_2,\widetilde Q_2')
 =-f_x(Q_1,Q_1')p-f_y(Q_1,Q_1')p'+o(\ell).
\]
Indeed the quadratic error is $O(\ell^2)$ and the phase remainder
is $O(M^{1-3c})$. Moreover
$(Q_1,Q_1')=(z_1+b_1h,b_1)+O(M^{1-2c})$; replacing the derivatives
of $f$ by their values at this affine pair produces only
$O(\ell M^{1-2c})=o(\ell)$. Similarly, the preceding $K^{-1}$ term
can be evaluated at the common affine pair with error $o(K^{-1})$.
We conclude that \cref{loc:cubic-relation} becomes
\begin{equation}
 P(h)=\frac{2p'}{b_1^3}\sec A
      -\frac{2\pi p}{b_1^2}\sin A\,\sec^2 A
      -\frac4{K b_1^2}\sec^2 A+o(\ell).
 \label{loc:opposite-expansion}
\end{equation}

Divide by $\ell$ and take a subsequence such that
$p/\ell\to\bar p$ coefficientwise and $K^{-1}/\ell\to\beta$.
At least one is nonzero, because
$\max(\norm{p/\ell}_{\rm coeff},K^{-1}/\ell)=1$.
Polynomial extraction now applies to $P/\ell$. Multiply its limiting
identity by $b^2\cos^2(2\pi(z+bh))$ and evaluate at every zero of
that cosine. One obtains
\begin{equation}
 2\pi\bar p(h)\sin(2\pi(z+bh))+4\beta=0
 \quad\hbox{at every cosine zero.}
 \label{loc:alternating-obstruction}
\end{equation}
The zeros form an unbounded arithmetic progression and their sine
values alternate between $1$ and $-1$. Thus $\bar p$ is bounded on
that progression and must be constant. The alternating signs then
force both that constant and $\beta$ to be zero, contradicting their
normalization. This excludes every positive limiting slope.

\subsection{A vanishing second slope}

It remains to treat $r\to0$. After a subsequence, the first term of
\cref{loc:cubic-relation} converges away from its poles to
$\sec(2\pi(z+bh))$, with $\abs{b}=1$. Write the second term as
\begin{equation}
 g(h)=\frac1{r^2S(h)},\qquad
 S(h)=\left(\frac{Q_2'(h)}r\right)^2
       \left(\cos(2\pi Q_2(h))+\frac{2-E_2(h)}K\right),
 \qquad w=\max(\abs{S(0)},r).
 \label{loc:slow-denominator}
\end{equation}
Uniformly on $[-1,1]$,
\begin{equation}
 \abs{S^{(m)}}\le C_m r^m\quad(1\le m\le4).
 \label{loc:slow-derivatives}
\end{equation}
To see every derivative involved, set $a=Q_2'/r$. Then $a$ is
bounded and bounded away from zero, and
\[
 \abs{a^{(m)}}\le C_m r^m M^{1-(m+1)c}\le C_m r^m
 \quad(1\le m\le4)
\]
by \cref{loc:scaled-phase-derivatives}. The derivatives through
order four of the cosine factor are $O(r^m)$ by the chain rule;
those of the shift are bounded by $C_mr^mM^{-c}/K$ by
\cref{loc:scaled-quotient}. The product rule proves
\cref{loc:slow-derivatives}. This is the point requiring the fifth
derivative of $Q_2$.

If $w\to0$, there is also the sharper fact
\begin{equation}
 \frac{\abs{S'(0)}}r\longrightarrow2\pi.
 \label{loc:slow-nonzero-slope}
\end{equation}
Indeed $a(0)^2=1$ and $\abs{Q_2'(0)}=r$. The cosine with its shift
equals $S(0)\to0$, so $\abs{\sin(2\pi Q_2(0))}\to1$.
Differentiating the factor $a^2$ contributes
$O(rwM^{1-2c})=o(r)$; differentiating the shift contributes
$O(rM^{-c}/K)=o(r)$. The remaining contribution is
$-2\pi\sin(2\pi Q_2(0))Q_2'(0)$, which proves the assertion.

Take a subsequence on which $r/w$ has a limit. If this limit is
positive, then $w\to0$, and Taylor's formula gives
\[
 \frac{S(h)}w=\frac{S(0)}w+\frac{S'(0)}w h+O(r^2/w).
\]
Here $r^2/w\le r\to0$. The coefficients are bounded; by
\cref{loc:slow-nonzero-slope}, a further subsequence gives a limiting
affine polynomial $A_0(h)$ with nonzero slope. Multiply
\cref{loc:cubic-relation} by $r^2w$. Its first term and its error
tend to zero locally away from the fast poles, while its second term
converges to $1/A_0(h)$ away from the one possible zero of $A_0$.
Polynomial extraction would therefore give a polynomial of degree
at most three equal to $1/A_0$. Multiplication by $A_0$ and evaluation
at its real zero is a contradiction.

We may consequently assume $r/w\to0$. Then $w=\abs{S(0)}$
eventually, and
\begin{equation}
 \frac{S(h)}{S(0)}=1+O(r/w)\longrightarrow1
 \quad\hbox{uniformly on }[-1,1].
 \label{loc:almost-constant-denominator}
\end{equation}
In particular $S$ has no zeros there for large $M$.
Let $T_3g$ be the cubic Taylor polynomial of $g$ at $0$.
If $w$ stays bounded away from zero, the reciprocal derivative formula
and \cref{loc:slow-derivatives} give $g^{(4)}=O(r^2)$ uniformly.
Thus $g-T_3g=O(r^2)$. Subtracting $T_3g$ from
\cref{loc:cubic-relation} leaves only the fast secant in the limit.
It cannot agree with a cubic polynomial, by multiplication by its
cosine and evaluation at a zero.

The last case is $w\to0$ and $r/w\to0$. To track also the possibly
large remainder of $g$, differentiate explicitly:
\begin{equation}
 g^{(4)}=\frac1{r^2}\left(
 \frac{24(S')^4}{S^5}
 -\frac{36(S')^2S''}{S^4}
 +\frac{6(S'')^2+8S'S'''}{S^3}
 -\frac{S''''}{S^2}\right).
 \label{loc:reciprocal-fourth}
\end{equation}
By \cref{loc:slow-derivatives,loc:slow-nonzero-slope},
$S'(h)/S'(0)=1+O(r)$ uniformly. Together with
\cref{loc:almost-constant-denominator}, this makes the first term
in \cref{loc:reciprocal-fourth} equal to
\[
 (24+o(1))\frac{S'(0)^4}{r^2S(0)^5}.
\]
Relative to that term, the remaining three displayed contributions
are respectively $O(w)$, $O(w^2)$, and $O(w^3)$: for example the
first ratio is bounded by $C\abs{S}\abs{S''}/\abs{S'}^2\le Cw$.
All these comparisons are uniform on $[-1,1]$. Put
\[
 a_M=\frac{S'(0)^4}{r^2S(0)^5}.
\]
Taylor's formula with integral remainder, valid on both sides of
zero, now yields
\begin{equation}
 g(h)-T_3g(h)=a_Mh^4+o(\abs{a_M})
 \quad\hbox{uniformly on }[-1,1].
 \label{loc:quartic-remainder}
\end{equation}
If $\abs{a_M}\to\infty$, subtract $T_3g$ from
\cref{loc:cubic-relation} and divide by $a_M$. The fast term tends
to zero away from its poles and the original error also tends to
zero after this division. Polynomial extraction would identify a
polynomial of degree at most three with $h^4$, which is impossible.
Otherwise pass to a subsequence with $a_M\to a\in\R$.
Without division, the same subtraction gives a limiting identity
between a cubic polynomial and
\[
 \sec(2\pi(z+bh))+ah^4.
\]
Move the quartic term to the polynomial side and multiply by the
cosine. Evaluation at any of its zeros again gives a contradiction.
These alternatives exhaust the vanishing-slope case.

\subsection{Uniform neighborhoods of successful sets}

We have proved that every Hausdorff limit of nonempty closures of
successful sets along $M\to\infty$ has measure zero. This implies
the uniform assertion of \cref{can:local-scarcity}, including the
neighborhood width in its statement. Indeed, if the assertion failed
for some $\eta>0$, then for every integer $n$ one could choose
$M_n\ge n$ and a successful set $\mathcal S_n\subset[-1,1]$ such
that
\[
 \operatorname{Leb}_{\R}\{x:\dist(x,\mathcal S_n)<1/n\}>\eta.
\]
These sets are nonempty. Compactness of the nonempty compact subsets
of $[-1,1]$ for Hausdorff distance supplies a subsequence of their
closures with limit $F$. For every $\varepsilon>0$, eventually
\[
 \{x:\dist(x,\mathcal S_n)<1/n\}
 \subset\{x:\dist(x,F)<\varepsilon\}.
\]
Thus every $\varepsilon$-neighborhood of $F$ has measure at least
$\eta$. Continuity of Lebesgue measure from above for these bounded
neighborhoods implies $\operatorname{Leb}_{\R}(F)\ge\eta$,
contrary to the result just proved. The empty successful set satisfies
the claimed estimate automatically. This completes the proof of
\cref{can:local-scarcity}.
